\documentclass{article}
\usepackage{amsmath}

\title{lec513}
\author{Dylan Tang}
\date{May 2026}

\begin{document}

\maketitle

\section{Review of GD and Backpropagation}
\subsubsection{Gradient Descent}
Gradient descent is given by:
\[
    W^{(t+1)}
    =
    W^{(t)}
    -
    \eta \nabla_{W}\mathcal{L}\left(W^{(t)}\right)
\]
\subsubsection{Backward Propogation}
\begin{itemize}
\item For a hidden layer $\ell$, recall that
\[
    Z^{(\ell)} = W^{(\ell)}A^{(\ell - 1)} + b^{(\ell)},
    \qquad
    A^{(\ell)} = \sigma\left(Z^{(\ell)}\right).
\]
\end{itemize}
Using the chain rule, the gradient with respect to $W^{(\ell)}$ is
\[
    \frac{\partial \mathcal{L}}{\partial W^{(\ell)}}
    =
    \frac{\partial \mathcal{L}}{\partial A^{(\ell)}}
    \frac{\partial A^{(\ell)}}{\partial Z^{(\ell)}}
    \frac{\partial Z^{(\ell)}}{\partial W^{(\ell)}}.
\]
Then, 
\[
    \frac{\partial A^{(\ell)}}{\partial Z^{(\ell)}} = \sigma'(Z^{(\ell)})
\]
and
\[
    \frac{\partial{Z}^{(\ell)}}{\partial W^{(\ell)}} = (A^{(\ell-1)})^T
\]
which gives
\[
    \boxed{
    \frac{\partial \mathcal{L}}{\partial W^{(\ell)}}
    =
    \left[
    \frac{\partial \mathcal{L}}{\partial A^{(\ell)}}
    \odot
    \sigma'\left(Z^{(\ell)}\right)
    \right]
    \left(A^{(\ell - 1)}\right)^T
    }
\]
Note that
\[
    \mathcal{L} = f(Y, \hat{Y}) = f(Y, \sigma(W^{(L)}A^{(L-1)})+b^{(L)}) = f(Y, \sigma(W^{(L)}\sigma(W^{(L-1)}A^{(L-2)})+b^{(L-1)})+b^{(L)}) = \dots
\]
\section{Convolution}
\[
    X = [a,b,c,d,e],
    \qquad
    Y = [-1,1]
\]

\begin{itemize}
    \item Think of \(X\) as the input signal.
    \item Think of \(Y\) as the filter/kernel.
    \item We slide \(Y\) across \(X\), computing the dot product at each slide.
\end{itemize}

\[
\begin{array}{c|ccccc}
X & a & b & c & d & e \\
\hline
Y & -1 & 1 &   &   &
\end{array}
\]

\[
    (-1)a + (1)b = -a + b
\]

\[
\begin{array}{c|ccccc}
X & a & b & c & d & e \\
\hline
Y &   & -1 & 1 &   &
\end{array}
\]

\[
    (-1)b + (1)c = -b + c
\]

\[
\begin{array}{c|ccccc}
X & a & b & c & d & e \\
\hline
Y &   &   & -1 & 1 &
\end{array}
\]

\[
    (-1)c + (1)d = -c + d
\]

\[
\begin{array}{c|ccccc}
X & a & b & c & d & e \\
\hline
Y &   &   &   & -1 & 1
\end{array}
\]

\[
    (-1)d + (1)e = -d + e
\]

\[
    X * Y
    =
    [-a+b,\,-b+c,\,-c+d,\,-d+e]
\]

\begin{itemize}
    \item Convolution length is $n - k + 1$, where $len(X) = 5$, $len(Y) = 2$
    \item This filter computes local differences between neighboring entries.
\end{itemize}
Consider the following image:
\[
X= 
\begin{bmatrix}
    0 & 0 & 0 & 0 \\
    0 & 1 & 1 & 0 \\
    0 & 1 & 1 & 0 \\
    0 & 0 & 0 & 0
\end{bmatrix}
\]
and had the filter 
\[
    Y = \begin{bmatrix}
        1 & -1 \\
        -1 & 1
    \end{bmatrix}
\]
Then, 
\[
    X*Y = \begin{bmatrix}
        1 & 0 & -1 \\
        0 & 0 & 0 \\
        -1 & 0 & 1
    \end{bmatrix}
\]
For the example above,
\begin{itemize}
    \item $1$ means the filter found a matching pattern
    \item $-1$ means the filter found an opposite pattern
    \item $0$ means the filter did not match in either direction
\end{itemize}

A CNN learns the combination of filters at each hidden layer that results in the best prediction.

Consider
\[
    X = \begin{bmatrix}
        a & b & c & d & e
    \end{bmatrix}
\]
Then, let
\[
    m_1 = \begin{bmatrix}
        m_{11} & m_{12}
    \end{bmatrix}
\]
and
\[
    m_2 = \begin{bmatrix}
            m_{21} & m_{22}
        \end{bmatrix}
\]
Then, 
\[
    X * m_1 * m_2
    =
    \begin{bmatrix}
    (am_{11}+bm_{12})m_{21} + (bm_{11}+cm_{12})m_{22} \\[0.5em]
    (bm_{11}+cm_{12})m_{21} + (cm_{11}+dm_{12})m_{22} \\[0.5em]
    (cm_{11}+dm_{12})m_{21} + (dm_{11}+em_{12})m_{22}
    \end{bmatrix}.
\]
\begin{itemize}
    \item Each element in the final series of convolutions contains elements in $X, m_1, m_2 \implies$ convolution gives each element information about the entire network 
    \item So, convolution is a very fast way to diseeminate information to all nodes of a layer
    \item Cascades of convolutions are called a \textbf{receptive field}. The end result contains information of all parts of the cascade.
\end{itemize}
\end{document}
